\section{Limitations and future works}


Our tasks are multiple-choice examination items, not clinical encounters: they carry no history
taking, no test ordering, no longitudinal follow-up, and the exact-match endpoint rewards the
option-selection component of diagnosis only. The $T=0$ property that makes the design clean also
bounds what it licenses --- results here do not extrapolate to tool-using clinical agents, which
is precisely the regime the general-domain study covers.

All models are from a single vendor. The heterogeneity arm of the original design, which mixes
vendors within a capability tier, is here restricted to mixing separately trained OpenAI families;
genuine cross-vendor error decorrelation may be larger than what we measure.

Panels are nested by construction (\S2.2), so the reported curves estimate accuracy as a routed
panel is extended, not the expected accuracy of independently redrawn panels at each size. The
coupling reduces variance on the architecture contrasts we care about and does not bias
within-item paired tests, but a reader interested in the independent-panel estimand should read
the curves accordingly.

The specialty router ranks by relevance, so larger panels append progressively less relevant
specialists. ``More experts'' and ``weaker experts'' are therefore entangled in the specialty-role
condition; we separate them with the generic-internist control, in which every panelist is an
undifferentiated internist and panel size is the only thing that varies.

Difficulty strata are defined by the pass rate of mid-tier models, some of which also appear in
the evaluation ladder. This inflates the level of the hard-stratum results but not the
within-item slope across panel sizes, which is computed on identical items.

\subsection{Success per 1K tokens}

[RESULTS]
